% Preámbulo reconstruido a partir del PDF compilado (pdfTeX, TeX Live 2023, Latin Modern).
\usepackage[utf8]{inputenc}
\usepackage[T1]{fontenc}
\usepackage{lmodern}
\usepackage[spanish,es-minimal,es-lcroman]{babel}
\usepackage{amsmath,amssymb,amsthm,mathtools,bm}
\usepackage{geometry}
\geometry{a4paper,margin=2.6cm,headheight=14pt}
\usepackage{enumitem}
\usepackage{booktabs}
\usepackage{fancyhdr}
\usepackage{xcolor}
\usepackage{tikz}
\usetikzlibrary{arrows.meta,decorations.pathmorphing}
\usepackage[most]{tcolorbox}
\usepackage[colorlinks=true,linkcolor=blue!60!black,citecolor=blue!60!black,urlcolor=blue!60!black]{hyperref}

\addto\captionsspanish{\renewcommand{\contentsname}{Índice}}

% ---- Cabeceras ----
\pagestyle{fancy}
\fancyhf{}
\fancyhead[LE]{\small\leftmark}
\fancyhead[RO]{\small\rightmark}
\fancyfoot[C]{\thepage}
\renewcommand{\headrulewidth}{0.4pt}
\fancypagestyle{plain}{\fancyhf{}\fancyfoot[C]{\thepage}\renewcommand{\headrulewidth}{0pt}}

% ---- Teoremas ----
\theoremstyle{plain}
\newtheorem{teorema}{Teorema}[chapter]
\newtheorem{proposicion}[teorema]{Proposición}
\newtheorem{lema}[teorema]{Lema}
\newtheorem{corolario}[teorema]{Corolario}
\theoremstyle{definition}
\newtheorem{definicion}[teorema]{Definición}
\newtheorem{ejemplo}[teorema]{Ejemplo}
\newtheorem{ejercicio}[teorema]{Ejercicio}
\theoremstyle{remark}
\newtheorem{observacion}[teorema]{Observación}

\definecolor{teoremaborde}{RGB}{60,60,140}
\definecolor{teoremafondo}{RGB}{246,246,252}
\definecolor{defborde}{RGB}{160,160,160}
\definecolor{deffondo}{RGB}{245,245,245}
\definecolor{panoramaborde}{RGB}{0,100,0}
\definecolor{panoramafondo}{RGB}{245,252,245}
\definecolor{intuicionborde}{RGB}{150,75,0}
\definecolor{intuicionfondo}{RGB}{255,250,240}

\tcolorboxenvironment{teorema}{enhanced,breakable,colback=teoremafondo,colframe=teoremaborde,
  boxrule=0.6pt,arc=1pt,left=4pt,right=4pt,top=2pt,bottom=2pt,before skip=8pt,after skip=8pt}
\tcolorboxenvironment{definicion}{enhanced,breakable,colback=deffondo,colframe=defborde,
  boxrule=0.4pt,arc=1pt,left=4pt,right=4pt,top=2pt,bottom=2pt,before skip=8pt,after skip=8pt}

\newtcolorbox{panorama}{enhanced,colback=panoramafondo,colframe=panoramaborde,boxrule=0.6pt,arc=1pt,
  title={Panorama del capítulo},fonttitle=\bfseries\small,left=6pt,right=6pt,before skip=10pt,after skip=12pt}
\newtcolorbox{intuicion}{enhanced,breakable,colback=intuicionfondo,colframe=intuicionborde,boxrule=0.6pt,arc=1pt,
  title={Intuición},fonttitle=\bfseries\small,left=6pt,right=6pt,before skip=10pt,after skip=10pt}

% ---- Notación ----
\newcommand{\R}{\mathbb{R}}
\newcommand{\N}{\mathbb{N}}
\newcommand{\C}{\mathbb{C}}
\newcommand{\Z}{\mathbb{Z}}
\newcommand{\E}{\mathbb{E}}
\renewcommand{\P}{\mathbb{P}}
\newcommand{\Q}{\mathbb{Q}}
\newcommand{\F}{\mathcal{F}}
\newcommand{\G}{\mathcal{G}}
\newcommand{\Hf}{\mathfrak{H}}
\newcommand{\Hc}{\mathcal{H}}
\newcommand{\Sc}{\mathcal{S}}
\newcommand{\D}{\mathbb{D}}
\newcommand{\ind}{\mathbf{1}}
\newcommand{\dd}{\mathrm{d}}
\newcommand{\Dom}{\operatorname{Dom}}
\newcommand{\Var}{\operatorname{Var}}
\newcommand{\Cov}{\operatorname{Cov}}
\newcommand{\Ent}{\operatorname{Ent}}
\newcommand{\tr}{\operatorname{tr}}
\newcommand{\sgn}{\operatorname{sgn}}
\newcommand{\Lip}{\operatorname{Lip}}
\newcommand{\ip}[2]{\left\langle #1, #2\right\rangle}
\newcommand{\norm}[1]{\left\lVert #1\right\rVert}
\newcommand{\abs}[1]{\left\lvert #1\right\rvert}
\newcommand{\cond}[2]{\E\left[#1\,\middle|\,#2\right]}
\newcommand{\law}{\overset{d}{=}}
\newcommand{\tolaw}{\xrightarrow{\ \mathrm{ley}\ }}
